% =============================================================================
% Presentación de Apoyo Visual: USS SpiderBot
% Taller de Programación I, Universidad San Sebastián. Año 2026.
% Compilar con: xelatex presentacion_hall_spiderbot.tex
% =============================================================================
\documentclass[10pt, aspectratio=169]{beamer}

% ─── Temas y Estilos ─────────────────────────────────────────────────────────
\usetheme{Madrid}
\usecolortheme{whale}

% ─── Paquetes ────────────────────────────────────────────────────────────────
\usepackage{fontspec}
\usepackage[spanish,es-noshorthands,es-noquoting]{babel}
\usepackage{xcolor}
\usepackage{listings}
\usepackage{booktabs}
\usepackage{amsmath}

% ─── Configuración de Colores Personalizados ────────────────────────────────
\definecolor{ussblue}{RGB}{12, 35, 64}      % Azul Marino Oficial USS
\definecolor{ussgold}{RGB}{201, 151, 65}    % Dorado Accento
\definecolor{bgcode}{RGB}{245, 245, 242}

\setbeamercolor{palette primary}{bg=ussblue, fg=white}
\setbeamercolor{palette secondary}{bg=ussblue!85, fg=white}
\setbeamercolor{palette tertiary}{bg=ussblue!70, fg=white}
\setbeamercolor{palette quaternary}{bg=ussblue, fg=white}
\setbeamercolor{titlelike}{parent=palette primary, fg=ussblue}
\setbeamercolor{structure}{fg=ussblue}
\setbeamercolor{item projected}{bg=ussgold, fg=ussblue}

% ─── Configuración de listings ──────────────────────────────────────────────
\lstdefinestyle{tinystyle}{
    language=Python,
    backgroundcolor=\color{bgcode},
    basicstyle=\tiny\ttfamily,
    keywordstyle=\color{blue}\bfseries,
    stringstyle=\color{purple},
    commentstyle=\color{gray}\itshape,
    breaklines=true,
    frame=single,
    rulecolor=\color{lightgray},
}

\newcommand{\codigo}[1]{\texttt{#1}}

% ─── Información de Portada ──────────────────────────────────────────────────
\title[USS SpiderBot]{USS SpiderBot: Robot Cuadrúpedo Articulado}
\subtitle{Solemne 3 --- Taller de Programación I}
\author[Moisés Amundarain]{Moisés Amundarain y Equipo}
\institute[USS]{Ingeniería Civil Informática\\ Universidad San Sebastián}
\date[Junio 2026]{2026}

% ─── Cuerpo de la Presentación ───────────────────────────────────────────────
\begin{document}

% Slide 1: Portada
\begin{frame}[plain]
    \titlepage
\end{frame}

% Slide 2: Introducción
\begin{frame}{Introducción y Propósito}
    \begin{columns}
        \begin{column}{0.5\textwidth}
            \begin{block}{Definición de Proyecto}
                El \textbf{USS SpiderBot} es un robot caminador cuadrúpedo de \textbf{8 grados de libertad (8-DoF)} orientado al reconocimiento autónomo en zonas de catástrofe y escombros.
            \end{block}
            \begin{block}{Desafío Tecnológico}
                Sustituir ruedas por extremidades articuladas para sortear desniveles físicos de forma estable utilizando control de bucle cerrado.
            \end{block}
        \end{column}
        \begin{column}{0.5\textwidth}
            \begin{alertblock}{Características Principales}
                \begin{itemize}
                    \item Modelado paramétrico en \textbf{OpenSCAD}.
                    \item Firmware multitarea asíncrono en \textbf{MicroPython}.
                    \item Estabilización inercial activa con IMU \textbf{MPU6050}.
                    \item Dashboard web Glassmorphism embebido.
                \end{itemize}
            \end{alertblock}
        \end{column}
    \end{columns}
\end{frame}

% Slide 3: Diseño Mecánico
\begin{frame}{Diseño Mecánico y Estructural (OpenSCAD)}
    \begin{columns}
        \begin{column}{0.5\textwidth}
            \begin{block}{Chasis Doble Deck}
                \begin{itemize}
                    \item \textbf{Placa Inferior ($r=65\text{mm}$):} Aloja perimetralmente los brackets de cadera.
                    \item \textbf{Placa Superior ($r=65\text{mm}$):} Cuenta con cuna central para protoboard de 400 puntos y marcos protectores.
                \end{itemize}
            \end{block}
            \begin{block}{Configuración Cinemática}
                Juntas en paralelo horizontal (\textbf{Pitch-Pitch}) para maximizar torque vertical y evitar cizalle dinámico.
            \end{block}
        \end{column}
        \begin{column}{0.5\textwidth}
            \begin{alertblock}{Tolerancias de Precisión (FDM)}
                \begin{itemize}
                    \item \textbf{Cavidad Servo:} $+0.3\text{mm}$ de holgura nominal.
                    \item \textbf{Bolsillos de Doble Profundidad:} Dual-depth ($30\text{mm}$ abajo para cables, $25\text{mm}$ arriba).
                    \item \textbf{Single-Arm Horn Locks:} Ranuras chavetadas que encastran el horn estriado a presión.
                \end{itemize}
            \end{alertblock}
        \end{column}
    \end{columns}
\end{frame}

% Slide 4: Arquitectura Electrónica
\begin{frame}{Arquitectura Electrónica y Conexiones}
    \begin{columns}
        \begin{column}{0.45\textwidth}
            \begin{block}{Bus I2C Compartido}
                El microcontrolador \textbf{ESP32 DevKit V1} comanda la lógica I2C en paralelo hacia dos esclavos:
                \begin{enumerate}
                    \item \textbf{PCA9685 ($0x40$):} Driver PWM de 12-bit.
                    \item \textbf{MPU6050 ($0x68$):} IMU de 6 ejes.
                \end{enumerate}
            \end{block}
        \end{column}
        \begin{column}{0.55\textwidth}
            \begin{alertblock}{Aislamiento de Potencia (UBEC)}
                \begin{itemize}
                    \item Los 8 servos consumen picos de hasta $4\text{A}$ combinados.
                    \item \textbf{UBEC (5V/3A):} Regula y separa la línea de fuerza externa de la bornera del PCA9685.
                    \item Previene caídas de tensión y reinicios (*brownouts*) del procesador ESP32.
                    \item \textbf{GND Unificado:} Referencia de voltaje común.
                \end{itemize}
            \end{alertblock}
        \end{column}
    \end{columns}
\end{frame}

% Slide 5: Firmware Multitarea
\begin{frame}{Firmware Multitarea Asíncrono (\codigo{uasyncio})}
    \begin{columns}
        \begin{column}{0.5\textwidth}
            \begin{block}{Planificación Cooperativa}
                Sustitución de llamadas bloqueantes (\codigo{time.sleep\_ms()}) por corrutinas no bloqueantes concurrentes:
                \begin{itemize}
                    \item \codigo{sensor\_updater()}: Adquisición MPU6050 y Sonar a $20\text{ Hz}$.
                    \item \codigo{locomotion\_loop()}: Lazo de cinemática y control servo.
                    \item \codigo{start\_server\_task()}: Socket TCP HTTP del dashboard.
                \end{itemize}
            \end{block}
        \end{column}
        \begin{column}{0.5\textwidth}
            \begin{alertblock}{Estado Compartido y Seguridad}
                \begin{itemize}
                    \item \textbf{state.py:} Singleton global que comparte la telemetría y órdenes de control en tiempo real.
                    \item \textbf{Lazo de Seguridad:} Si el sonar detecta un obstáculo a menos de $15\text{ cm}$, se cancela la locomoción en menos de $25\text{ ms}$, regresando a pose segura con pitido acústico.
                \end{itemize}
            \end{alertblock}
        \end{column}
    \end{columns}
\end{frame}

% Slide 6: Marcha de Gateo
\begin{frame}{Algoritmo de Locomoción: Marcha de Gateo (Crawl Gait)}
    \begin{columns}
        \begin{column}{0.5\textwidth}
            \begin{block}{Polígono de Sustentación}
                Estabilidad estática garantizada: en fase de oscilación (swing), el centro de masa (CoM) siempre reside dentro del triángulo formado por los tres puntos de apoyo en el suelo.
            \end{block}
            \begin{block}{Ciclo Cinemático (5 Fases)}
                Secuencia repetitiva: Stance Shift general $\rightarrow$ Paso Pata FR $\rightarrow$ Paso Pata RR $\rightarrow$ Paso Pata FL $\rightarrow$ Paso Pata RL.
            \end{block}
        \end{column}
        \begin{column}{0.5\textwidth}
            \begin{alertblock}{Valores Angulares del Ciclo}
                \begin{itemize}
                    \item \textbf{Coxa (Cadera):} Avance derecho $110^\circ$, retroceso $70^\circ$. Espejado en el lado izquierdo (avance $75^\circ$, retroceso $105^\circ$).
                    \item \textbf{Fémur (Rodilla):} Apoyo en el suelo a $60^\circ / 120^\circ$. Oscilación (swing) a $90^\circ$ para levantar la extremidad.
                \end{itemize}
            \end{alertblock}
        \end{column}
    \end{columns}
\end{frame}

% Slide 7: Estabilización Activa
\begin{frame}{Lazo Cerrado de Estabilización Inercial Activa}
    \begin{columns}
        \begin{column}{0.5\textwidth}
            \begin{block}{Cálculo de Inclinación}
                La IMU MPU6050 mide los ángulos de Pitch ($\theta_p$) y Roll ($\theta_r$) a una tasa de $20\text{ Hz}$.
            \end{block}
            \begin{block}{Compensación Proporcional}
                Inyección de offsets angulares dinámicos en los servos de rodilla (fémures) de las patas en contacto:
                $$\theta_{compensado} = \theta_{base} \pm (K_p \cdot \theta_{sensor})$$
                Donde el coeficiente $K_p = 1.2$ compensa proporcionalmente la caída del chasis.
            \end{block}
        \end{column}
        \begin{column}{0.5\textwidth}
            \begin{alertblock}{Comportamiento Reactivo}
                \begin{itemize}
                    \item Las patas del lado inclinado hacia abajo se estiran de forma automática.
                    \item Las patas del lado elevado se contraen.
                    \item El plano de la base superior se mantiene horizontal en tiempo real.
                    \item Las patas en fase de aire (swing) ignoran la corrección para no perder avance.
                \end{itemize}
            \end{alertblock}
        \end{column}
    \end{columns}
\end{frame}

% Slide 8: Servidor y Dashboard
\begin{frame}{Dashboard Web Embebido de Telemetría y Mando}
    \begin{columns}
        \begin{column}{0.5\textwidth}
            \begin{block}{Dashboard Autocontenido (Glassmorphism)}
                \begin{itemize}
                    \item Archivo HTML, CSS y JS unificado de $24\text{ KB}$ servido de forma asíncrona.
                    \item \textbf{100\% Offline:} Sin CDNs o llamadas de red externas para evitar dependencias locales.
                    \item Enrutamiento por bloques de 512 bytes para evitar desbordamientos en la memoria de la ESP32.
                \end{itemize}
            \end{block}
        \end{column}
        \begin{column}{0.5\textwidth}
            \begin{alertblock}{Funciones de la Interfaz}
                \begin{itemize}
                    \item \textbf{Mando por Teclado:} Soporte de teclas físicas (W-A-S-D), barra espaciadora (parada) y tecla R (reposo).
                    \item \textbf{Graficador SVG:} Visualizador vectorial 2D en tiempo real del plano inclinado del chasis inercial.
                    \item \textbf{Telemetría en Vivo:} Lectura constante de distancia e inclinación.
                \end{itemize}
            \end{alertblock}
        \end{column}
    \end{columns}
\end{frame}

% Slide 9: Simulación y Fabricación
\begin{frame}{Validación en Wokwi e Impresión 3D}
    \begin{columns}
        \begin{column}{0.5\textwidth}
            \begin{block}{Simulador Virtual Wokwi}
                \begin{itemize}
                    \item Cableado digital completo de 8 servos y sensores en \codigo{diagram.json}.
                    \item \textbf{Control Fallback Directo:} En ausencia del expansor PCA9685, el firmware conmuta a modulación PWM nativa de la ESP32.
                    \item Conexión a través del canal virtual `Wokwi-GUEST` para interactuar con el servidor real.
                \end{itemize}
            \end{block}
        \end{column}
        \begin{column}{0.5\textwidth}
            \begin{alertblock}{Especificaciones de Slicing (Creality Print)}
                \begin{itemize}
                    \item \textbf{Altura de Capa:} $0.2\text{ mm}$ con PLA/PETG.
                    \item \textbf{Líneas de Pared:} Mínimo 4 perímetros (esencial para la fuerza de los agujeros roscados M2).
                    \item \textbf{Relleno:} $30\%$ en patrón \textbf{Gyroid} (giroidal) por su rigidez mecánica tridimensional isotrópica.
                    \item Proyecto de plato listo guardado en [spiderbot.3mf].
                \end{itemize}
            \end{alertblock}
        \end{column}
    \end{columns}
\end{frame}

% Slide 10: Conclusiones
\begin{frame}{Conclusiones y Trabajo Futuro}
    \begin{block}{Logros Obtenidos}
        \begin{itemize}
            \item Integración exitosa de modelado CAD, electrónica embebida, control asíncrono y desarrollo de red en la ESP32.
            \item Estabilización proporcional inercial y evasión de obstáculos dinámicas completamente funcionales.
            \item Estructura paramétrica en OpenSCAD lista para laminación FDM.
        \end{itemize}
    \end{block}
    \begin{alertblock}{Mejoras Propuestas (Trabajo Futuro)}
        \begin{itemize}
            \item \textbf{Cinemática Inversa (IK):} Control completo de 3-DoF por pata para trayectorias de pie exactas.
            \item \textbf{Visión Inteligente:} Integrar cámara web ESP32-CAM y algoritmos de detección local para rescates.
            \item \textbf{Marcha Dinámica:} Programar trote alternado con fases de vuelo.
        \end{itemize}
    \end{alertblock}
\end{frame}

\end{document}
